\acterecto{Devoir envers l’Église et ses ministres}
\usage{quot.}
\ouverture{P}{our}{m’acquitter} devant vous, ô mon Dieu, de mon devoir de chrétien,

je me propose aujourd’hui d’offrir toutes mes prières, supplications, sacrifices et œuvres à l’intention de la sainte Église catholique,

pour son triomphe sur les erreurs, son affermissement dans la vérité,

et surtout pour le retour d’un saint pape, digne Vicaire de Jésus-Christ, selon la promesse faite à saint Pierre.

Je m’unis aux intentions de l’Église militante, vous priant instamment pour la sanctification, la piété, le zèle, la persévérance et la subsistance temporelle de nos saints prêtres, vos fervents ministres, desquels nous recevons la suavité de vos sacrements si désirables, et de leurs instituts qui assurent leur rigoureuse formation. Je désire aussi que mes prières soient portées au soulagement des âmes des religieux retenues en purgatoire.

\invocation{Saint Pierre, prince des Apôtres, priez pour nous.\par
Saint Jean, apôtre bien-aimé, priez pour nous.\par
Saint Pie X, défenseur de la foi, priez pour la sainte Église.}
\culdelampe

\acte{Devoir envers la famille}
\usage{quot.}
\ouverture{P}{our}{m’acquitter} de mes devoirs filiaux, honorer mes parents et me rendre digne de ma lignée, je veux prier aujourd’hui pour tous les membres de ma famille,

pour leur conversion, leur conservation dans le bien, leur soulagement dans les épreuves temporelles, leur persévérance dans la grâce, et leur sanctification selon les desseins et dans les voies de Dieu.

Je souhaite grandement confier le repos des âmes de mes ancêtres et de tous mes défunts, à votre infinie miséricorde, assuré ainsi qu’aucun ne demeure oublié.

\invocation{Jésus, Marie, Joseph, éclairez-nous, secourez-nous, sauvez-nous.\par
Saint Joseph, chef de la sainte Famille, gardez nos foyers.}
\culdelampe

\acte{Devoir envers la patrie}
\usage{quot.}
\ouverture{J}{’ose}{élargir encore} mes intentions, Seigneur, pour m’acquitter de mes devoirs envers ma patrie terrestre.

Que mes œuvres soient retenues par votre souveraine libéralité, ô mon Dieu, pour le salut de la France, la conversion de ses chefs, le retour de ses lois à la loi du Christ et à son saint règne social ; la fin du joug antichristique dans ses institutions ; et le renouvellement des promesses de son baptême.

J’offre aussi ma prière pour le repos des âmes de mes compatriotes morts dans l’amitié de Dieu, par la guerre ou par l’oppression.

\invocation{Sacré-Cœur de Jésus, sauvez la France.\par
Notre-Dame de l’Assomption, patronne principale de la France, priez pour nous.\par
Notre-Dame de France, priez pour nous.\par
Saint Michel Archange, protecteur de la France, priez pour nous.\par
Saint Remi, apôtre des Francs, saint Denis, saint Martin, apôtres des Gaules, priez pour nous.\par
Saint Louis, roi de France, priez pour nous.\par
Sainte Clotilde, sainte Geneviève, priez pour nous.\par
Sainte Jeanne d’Arc, libératrice de la patrie, priez pour nous.\par
Sainte Thérèse de l’Enfant-Jésus, patronne secondaire de la France, priez pour nous.\par
Saints martyrs de la Révolution, priez pour la France.}
\invocation{Vive le Christ qui aime les Francs !}
\culdelampe

\acte{Devoir envers les pauvres, les persécutés, les affligés}
\usage{quot.}
\ouverture{J}{e}{ne veux pas} passer ce jour, Seigneur, sans me souvenir des pauvres, des persécutés, des affligés, des solitaires, de tous ceux qui souffrent dans leur corps ou leur âme, et que personne ne recommande à Dieu.

Que mes prières montent pour eux ; surtout pour ceux qui souffrent au nom de Notre-Seigneur Jésus-Christ, afin que la justice divine soit fléchie et que la miséricorde les visite.

\invocation{Sainte Marie, salut des infirmes, consolatrice des affligés, refuge des pécheurs, priez pour eux.\par
Saint Vincent de Paul, père des pauvres, priez pour nous.}
\culdelampe

\acte{Devoir envers moi-même}
\usage{quot.}
\ouverture{A}{ussi,}{Seigneur,} car charité bien ordonnée ne saurait s’oublier, je me propose, en vue de mon propre salut, de vous offrir aujourd’hui ma vigilance, mes combats, mes pensées et mes œuvres, vous priant pour le salut de mon âme.

Que par votre grâce, tout en moi soit purifié de ses attaches, ordonné à votre volonté, et sanctifié selon vos desseins.

Qu’aucune vile illusion ne voile à mes yeux l’exigence de la sainteté.

Faites, Seigneur, que ce jour, en dépit de mes misères, me conduise plus près de vous.

\invocation{Saint Gabriel, et vous, mon bon Ange gardien, priez pour moi.}